% !TEX root = ../main_mk.tex
\chapter{Преглед на областа}\label{ch:background}

\section{Пренос на стил на текст}

Преносот на стил на текст го препишува текстот така што се менува едно
негово својство, стилот, а содржината останува иста \citep{jin2022survey}.
Што се смета за стил зависи од задачата: сентимент, формалност,
учтивост или начинот на пишување на одреден автор. Последниот вид е
најстариот пример во областа: \citet{xu2012shakespeare} го парафразираа
современиот англиски јазик во Шекспировиот и обратно, користејќи ги
паралелните современи верзии на неговите драми.

\paragraph{Надгледуван пренос.} Кога постојат паралелни податоци, парови
текстови со иста содржина во два стила, задачата е обично учење од
секвенца во секвенца. Најпознатото такво множество, GYAFC, има околу
110.000 парови неформални и формални реченици, напишани од луѓе ангажирани
преку платформи за краудсорсинг \citep{rao2018gyafc}. Паралелни податоци
за стилот на одреден автор практично не постојат: никој не ја пишува истата
песна двапати, со гласовите на двајца поети.

\paragraph{Ненадгледуван пренос.} Затоа поголемиот дел од литературата
работи со непаралелни податоци, текстови во секој стил што не се преводи
еден на друг. Едно такво истражување ги одделува содржината и стилот во
латентна репрезентација и ја декодира содржината со друг стил
\citep{shen2017crossalign}. Друго пак директно го уредува текстот: ги брише
зборовите што го означуваат изворниот стил и ги пополнува празнините со
зборови типични за целниот стил \citep{li2018delete}. Трет труд го третира
преносот како парафразирање: модел за парафразирање најпрвин го
отстранува стилот, а модел обучен врз текстови од еден стил потоа ја
препишува неутралната парафраза во тој стил \citep{krishna2020strap}. И
тука е применето учење со поттикнување, со награди за целниот стил и за
зачувување на содржината \citep{luo2019dualrl}.

\paragraph{Барања до големи модели.} Големите јазични модели можат да
вршат пренос на стил без никаква обука, само со инструкција и неколку
примери во барањето (prompt) \citep{reif2022recipe}. Методите со
пополнување (infilling) го постигнуваат уредувањето на текстот со праќање барања до моделот: зборовите
што го носат стилот се означуваат или маскираат, а јазичниот модел ги
избира нивните замени, така што остатокот од текстот, а со него и
содржината, останува непроменет \citep{pegf}. Вториот од претходните
пристапи во оваа работа (Поглавје~\ref{ch:path}) го следи овој дизајн.

\paragraph{Поезија и авторски стил.} Поезијата додава ограничувања што
прозата ги нема: стихови, строфи, метрика и рима.
\citet{tikhonov2018guesswho} генерираат поезија во стилот на одредени
автори, на англиски и на руски јазик, и ги оценуваат резултатите и
автоматски и со човечки читатели. Нивните модели генерираат нови песни
од корпусот на авторот; задачата во оваа работа е во еден поглед
потешка, бидејќи генерираната песна мора да ја задржи и содржината на
дадена изворна песна, напишана од друг автор.

\section{Стилометрија и утврдување авторство}

Стилометријата го мери стилот преку пребројливи карактеристики на
текстот, а нејзината главна примена е утврдување авторство:
одлучување кој го напишал дадениот текст \citep{stamatatos2009survey}.
Класичната студија на спорните текстови \emph{Federalist} покажа дека
честотите на чести функциски зборови, кои авторот ги користи без да
размислува за нив, го откриваат авторот подобро отколку содржината
\citep{mosteller1964}. Мерката Delta на Burrows \citep{burrows2002delta},
растојание врз стандардизирани честоти на најчестите зборови, е стандардна основа. 
Современото утврдување авторство комбинира
многу видови карактеристики ($n$-грами од зборови и знаци, шеми на
зборовни видови, лексичко богатство, должина на реченици и стихови) со
надгледуван класификатор \citep{stamatatos2009survey}. Лексичкото
богатство зависи од должината на текстот, па за да се споредуваат кратки
и долги песни потребна е мерка што не расте со текстот; подвижниот
просечен однос на типови и појавници (MATTR) го пресметува односот во
лизгачки прозорец со фиксна големина \citep{covington2010mattr}.

Оваа работа ја користи стилометријата во две улоги. Профилот на
карактеристиките на секој автор дефинира што бара наградата, а
класификатор на автори мери дали генерираната песна звучи како
песна на тој автор. Кога истиот класификатор и добива награда, и се користи
за мерење на успехот, моделот може да научи да го задоволи
класификаторот без да пишува како авторот. Затоа евалуацијата користи втор
класификатор што обуката никогаш не го гледа
(Поглавје~\ref{ch:stylometry}).

\section{Дестилација на знаење и адаптација со низок ранг}

Дестилацијата на знаење обучува мал модел, \emph{ученик}, да имитира
поголем модел, \emph{учител} \citep{hinton2015distilling}. За генерирање
текст наједноставната и најефикасната форма го обучува ученикот врз
излезите на учителот, како да се референтни податоци (дестилација на ниво
на секвенца, \citealp{kim2016seqkd}). Големите јазични модели го практикуваат
ова како вообичаен начин за градење податоци за обука: силен модел пишува
примери за задачата, а помал модел се фино подесува врз нив
\citep{wang2023selfinstruct}. Учениците во оваа работа се обучуваат на
истиот начин, врз песни напишани од учителот.

Финото подесување на сите тежини на модел со милијарди параметри бара
повеќе меморија отколку што има една графичка картичка. Адаптацијата со
низок ранг (LoRA, \citealp{hu2022lora}) ги замрзнува оригиналните тежини
и, за избрани тежински матрици, учи корекција во облик на производ на две
мали матрици. Параметрите што се обучуваат се дел од процентот на
моделот, а по обуката корекцијата може да се спои со тежините.

\section{Учење со поттикнување кај јазичните модели}

Надгледуваното фино подесување го учи моделот да имитира примери.
Учењето со поттикнување, пак, му дозволува на моделот да генерира, го
оценува секој излез со \emph{награда} и го менува моделот така што
излезите со висока награда стануваат поверојатни. Кај јазичните модели
наградата најпрвин се учеше од човечки преференции меѓу парови излези
\citep{christiano2017deep, ziegler2019finetuning}; ова учење со
поттикнување од човечка повратна информација (RLHF) е начинот на кој се
обучуваат агентите кои следат инструкции \citep{ouyang2022instructgpt}.
Алгоритмот за оптимизација во тој труд е Proximal Policy Optimization
(PPO, \citealp{schulman2017ppo}), кој, покрај моделот што се обучува, користи
и втора мрежа со слична големина за проценка на вредноста на секоја
состојба.

Group Relative Policy Optimization (GRPO, \citealp{shao2024deepseekmath})
ја отстранува таа втора мрежа. За секое барање генерира група излези, ги
оценува, и како предност (advantage) ја користи наградата на секој излез
во однос на средната вредност и стандардната девијација на неговата група.
Казна за отстапувањето од референтниот модел го држи обучуваниот модел
блиску до почетната точка. GRPO е воведен за математичко расудување, а
подоцна се користи за обука на модели за расудување со награди што можат
автоматски да се проверат \citep{deepseek2025r1}. Тука наградата се
пресметува, а не се учи од преференции: тоа е стилометрискиот профил на
целниот автор заедно со мерка за зачуваноста на содржината.

Секоја награда е само приближување до она што се бара, а модел што доволно
силно ја оптимизира наоѓа начини да добие висока оценка без да ја изврши
задачата. Оваа \emph{измама на наградата} (reward hacking,
\citealp{skalse2022hacking}) расте со јачината на оптимизацијата
\citep{gao2023overoptimization}, и затоа наградата и евалуацијата во оваа
работа користат различни мерења. Обуката ја користи библиотеката TRL
\citep{vonwerra2020trl}.

\section{Обработка на македонскиот јазик}

Македонскиот е јужнословенски јазик што се пишува на кирилица, со богата
глаголска морфологија, постпозитивен определен член со три форми и речиси
без падежи кај именките. Има многу помалку ресурси од англискиот или од
поголемите словенски јазици. Проектот CLASSLA нуди означување и
лематизација за јужнословенските јазици, изградени врз алатката Stanza
\citep{qi2020stanza, ljubesic2019neural, tercon2023classla}; неговите
модели за македонски доделуваат лема на секој збор, универзален
облик на зборот и морфосинтаксички опис според MULTEXT-East
\citep{erjavec2012multext}. Големите општи модели пишуваат на македонски,
но со отстапки во граматиката и во усогласувањето на зборовите
\citep{arizankovska2024chatgpt}, а малите отворени модели што се користат
како ученици во оваа работа се обучени врз многу малку текст на
македонски.
Првиот отворен основен модел изграден специјално за македонски, со модел
од 8 милијарди параметри, корпус од 3,5 милијарди зборови, множество
инструкции и мерно множество од седум задачи, се појавува дури во 2025
година \citep{krsteski2025macedonian}.
